\section{离开 OpenAI 与人才哲学}

\subsection{创立 Anthropic 的真正原因}

Dario 在 OpenAI 工作约 5 年，最后两年担任研究副总裁。他和 Ilya Sutskever 在 2016-2017 年共同制定了研究方向，Dario 在 Scaling 方面全力推进：GPT-2、GPT-3、RLHF、辩论、放大、可解释性。

关于离开原因，他澄清了两个常见误解：
\begin{itemize}
    \item \textbf{不是}因为微软投资（这是错误说法）
    \item \textbf{不是}因为商业化（他本人构建了 GPT-3 并推动了其商业化）
\end{itemize}

真正原因是对``如何做''有不同愿景——更谨慎、更直接、更诚实、建立信任。\textit{``安全如何才能不仅仅是我们为了招聘而说的话？''}

\begin{importantbox}{``去做你的愿景''}
\textit{``在别人的组织里争论自己的愿景是极其低效的……更有成效的做法是离开，去做一个干净的实验。''}

这个哲学也体现在 Race to the Top 策略中：\textit{``最终谁赢并不重要，只要每个人都在互相学习对方的好做法。''} Race to the Bottom 才是真正的敌人。

Anthropic 是\textit{``一群不完美的人，不完美地瞄准某个永远无法完美实现的理想。''}
\end{importantbox}

\subsection{人才密度 vs 人才数量}

\textit{``这个观点每个月都更加正确。''}

Dario 的思想实验：100 个超级聪明且方向一致的人 $>$ 1000 个其中 200 个优秀 + 800 个普通大厂员工。原因：

\begin{itemize}
    \item \textbf{高密度} = 信任、互相激励、所有人都看到更大目标
    \item \textbf{低密度} = 需要流程、护栏、政治斗争、各自为政
\end{itemize}

Anthropic 从 300 人快速扩张到 800 人，然后从 800 到约 950 人明显放慢。Dario 认为在 1000 人左右有一个``拐点''——需要更加谨慎。他招了很多理论物理学家，因为\textit{``他们学东西非常快。''}

引用 Steve Jobs：\textit{``A 类人才希望环顾四周看到其他 A 类人才。''} 看到不是全力以赴追求使命的人\textit{``是令人沮丧的。''}

\subsection{什么造就优秀的 AI 研究者}

第一品质：\textbf{开放心态}。Dario 自己的经历——看到了和所有人一样的数据，不是更好的程序员，只是\textit{``愿意用新的眼光看事物。''}

\textit{``Scaling Hypothesis 的发现是简单和愚蠢的。任何人都可以做到。''} 但事实上只有``个位数的人''推动了整个领域的这个认知。

\begin{knowledgebox}{给年轻人的建议}
\begin{itemize}
    \item 直接开始玩模型——\textit{``这些模型是没人真正理解的新产物''}
    \item 探索未充分研究的领域：可解释性（只有约 100 人 vs 架构方向的 10000 人）、长期学习、评估、多 Agent
    \item \textit{``滑向冰球要去的地方''}
    \item 经验有时是劣势——开放心态往往来自对领域的新鲜感
\end{itemize}
\end{knowledgebox}

\subsection{本章小结}

Anthropic 的创立源于``去做你自己的愿景''的哲学。人才密度是核心竞争力——高密度团队无需繁重流程就能高效运转。优秀 AI 研究者的核心品质是开放心态，而非技术经验。

